% ---------------------------------------------------------------------------
%  Section 4 -- Experiments
%
%  4.1-4.3 are complete: every hyperparameter below was read from the config
%  embedded in the trained checkpoints, not from memory.
%
%  4.4 is a skeleton. Numbers are marked "--" and each table carries a comment
%  saying which file the values come from.
%
%  CASE COUNTS VERIFIED: 1251 + 1350 = 2601 studies; a 70/20/10 split gives
%  1821 / 520 / 260. Cross-checked against the training log -- 1821 train
%  cases at batch 32 with drop_last is 56 steps/epoch, and 200 epochs x 56 =
%  11200 steps, exactly what EXPERIMENTS.md recorded for the ImageVAE run.
%
%  STILL TO CONFIRM:
%    - The checkpoint available here was trained with subregion_mode=True,
%      i.e. a four-channel annotation autoencoder. The text below describes
%      the binary configuration. Re-check 4.2 against the config of the run
%      that actually produces the reported numbers.
%    - Training hardware (marked TODO inline).
% ---------------------------------------------------------------------------

\section{Experiments}
\label{sec:experiments}

% ---------------------------------------------------------------------------
\subsection{Dataset and Preprocessing}
\label{subsec:dataset}

We evaluate on the BraTS 2023 and 2024 adult glioma
collections~\cite{brats_menze,brats_bakas,brats2021}, which provide
$1{,}251$ and $1{,}350$ studies with publicly released annotations, giving
$2{,}601$ in total. Each study provides four co-registered, skull-stripped
sequences (T1, T1CE, T2 and FLAIR) at
$240 \times 240 \times 155$ with a voxel-wise expert annotation. As described
in Section~\ref{subsec:problem}, we take the whole tumour as the binary
target.

Intensities are normalised per sequence and per study by $z$-scoring over
brain voxels only, so that background air does not bias the statistics and
scanner-dependent ranges are placed on a common scale. Each study is then
cropped to a fixed $128^{3}$ window centred on the centroid of the annotation
bounding box. A fixed crop allows annotation latents to be computed once and
cached, and this window retains the complete tumour for $97\%$ of studies.

Studies are partitioned patient-wise into training, validation and test sets
in a $70{:}20{:}10$ ratio with a fixed seed, giving $1{,}821$, $520$ and
$260$ studies respectively. All configurations of a given study fall in the
same partition, so no case is ever evaluated under one subset of sequences
while having been trained on under another.

% ---------------------------------------------------------------------------
\subsection{Implementation Details}
\label{subsec:implementation}

\paragraph{Autoencoders.}
Both autoencoders use a residual convolutional encoder--decoder with two
residual units per level and two downsampling stages, mapping $128^{3}$ inputs
to $32^{3}$ latents. The image autoencoder uses channel widths
$(64, 128, 256)$ and a KL weight of $\beta_{X} = 10^{-4}$; the annotation
autoencoder uses $(32, 64, 128)$, a latent width of $C_{S} = 4$ and
$\beta_{S} = 10^{-2}$. The larger KL weight on the annotation branch reflects its
simpler target, since a binary mask needs less latent capacity than a
four-sequence volume.

Both are trained for 200 epochs with AdamW~\cite{adam}, learning rate
$10^{-4}$, weight decay $10^{-4}$ and batch size 32, with early stopping after
25 validation checks without improvement. Modality dropout is applied when
training the image autoencoder as well as the diffusion model, so the encoder learns to
represent partial inputs rather than meeting them for the first time at
inference. Both autoencoders are frozen thereafter.

\paragraph{Diffusion model.}
The denoising network is a 3D U-Net with base width 64, channel multipliers
$(1, 2, 4)$ and two residual blocks per level, totalling $31.3$M parameters.
The image latent is projected from 256 to 64 channels before being
concatenated with the noisy annotation latent, so that the conditioning does not dominate the input by width. Training uses $T = 1000$ timesteps with a
linear noise schedule, $2 \times 10^{5}$ optimisation steps, batch size 64 and
AdamW with learning rate $10^{-4}$ and weight decay $10^{-4}$. The modality
vector $\mathbf{m}$ is resampled independently for every example, uniformly
over the 15 non-empty configurations.

\paragraph{Inference.}
Sampling uses DDIM~\cite{ddim} with 50 steps. We draw $N = 5$ samples per
case, average them to obtain the prediction, and threshold at $\tau = 0.5$.
Figure~\ref{fig:n-sweep} shows how whole-tumour Dice varies with $N$ on the
full-modality configuration. Accuracy rises steeply for the first few samples
and then flattens, so $N = 5$ sits at the point where further samples cost
inference time without improving the ensemble.

\begin{figure}[t]
    \centering
    \includegraphics[width=0.62\textwidth]{n_sweep.pdf}
    \caption{Whole-tumour Dice against the number of sampling seeds $N$, with
    all four sequences available. The curve flattens after five samples, which
    is the value used throughout.}
    \label{fig:n-sweep}
\end{figure} All experiments use a fixed random seed and are implemented in
PyTorch; training was performed on NVIDIA H200 GPUs. % TODO: confirm hardware

% ---------------------------------------------------------------------------
\subsection{Evaluation Protocol and Metrics}
\label{subsec:metrics}

Every test case is evaluated under all 15 modality configurations, giving a
complete case $\times$ configuration grid. Segmentation quality is reported
with the Dice coefficient and Intersection over Union between the thresholded
ensemble and the annotation. Predictive uncertainty is reported with
$U(X_{\mathbf{m}})$ from Eq.~\eqref{eq:uncertainty}, the mean sample variance
restricted to the region where samples can disagree, and with the paired
difference $\Delta U_{i}(\mathbf{m})$ of Eq.~\eqref{eq:delta-u} relative to
the fully observed input for the same case.

Because the autoencoder is frozen, its reconstruction fidelity upper-bounds
what any diffusion model operating in that latent space can achieve. We therefore report autoencoder
reconstruction separately in Table~\ref{tab:autoencoder}, so that the
segmentation results in Table~\ref{tab:modality} can be read against that
ceiling.

% ---------------------------------------------------------------------------
\subsection{Results}
\label{subsec:results}

% --- Table 1: autoencoder ceiling -----------------------------------------
% SOURCE: final validation row of the image_vae and mask_vae run logs
%         (logs/runs/<run_id>/val.csv)
\begin{table}[t]
    \centering
    \caption{Reconstruction fidelity of the frozen autoencoders. The
    annotation autoencoder bounds the segmentation performance attainable in
    latent space.}
    \label{tab:autoencoder}
    \begin{tabular}{lcccc}
        \toprule
        Autoencoder & Dice & IoU & SSIM & PSNR \\
        \midrule
        Image ($128^{3} \to 32^{3}$)      & --   & --   & --   & --   \\
        Annotation ($128^{3} \to 32^{3}$) & --   & --   & --   & --   \\
        \bottomrule
    \end{tabular}
\end{table}

% --- Table 2: the main result ---------------------------------------------
% SOURCE: eval_output/<run>/summary.csv, one row per combo.
% Bit order is (FLAIR, T1ce, T1, T2) -- keep consistent with Sec. 3.1.
% Delta U is computed against the 1111 row of the same case, then averaged.
\begin{table}[t]
    \centering
    \caption{Segmentation accuracy and predictive uncertainty across all 15
    modality configurations, grouped by the number of available sequences.
    $\Delta U$ is the paired increase in uncertainty relative to the fully
    observed input for the same case.}
    \label{tab:modality}
    \begin{tabular}{llcccc}
        \toprule
        $|\mathbf{m}|$ & Available sequences & Dice & IoU & $U$ & $\Delta U$ \\
        \midrule
        4 & FLAIR + T1ce + T1 + T2 & -- & -- & -- & --   \\
        \midrule
        \multirow{4}{*}{3}
          & FLAIR + T1ce + T1      & -- & -- & -- & -- \\
          & FLAIR + T1ce + T2      & -- & -- & -- & -- \\
          & FLAIR + T1 + T2        & -- & -- & -- & -- \\
          & T1ce + T1 + T2         & -- & -- & -- & -- \\
        \midrule
        \multirow{6}{*}{2}
          & FLAIR + T1ce           & -- & -- & -- & -- \\
          & FLAIR + T1             & -- & -- & -- & -- \\
          & FLAIR + T2             & -- & -- & -- & -- \\
          & T1ce + T1              & -- & -- & -- & -- \\
          & T1ce + T2              & -- & -- & -- & -- \\
          & T1 + T2                & -- & -- & -- & -- \\
        \midrule
        \multirow{4}{*}{1}
          & FLAIR                  & -- & -- & -- & -- \\
          & T1ce                   & -- & -- & -- & -- \\
          & T1                     & -- & -- & -- & -- \\
          & T2                     & -- & -- & -- & -- \\
        \bottomrule
    \end{tabular}
\end{table}

% --- Figure: the central claim, one plot ----------------------------------
% Uncertainty against number of available sequences. Box plot over cases, or
% mean with standard error. This figure is the paper's main evidence.
\begin{figure}[t]
    \centering
    \includegraphics[width=0.75\textwidth]{uncertainty_vs_modalities.pdf}
    \caption{Predictive uncertainty against the number of available sequences.
    Each case contributes one observation per configuration, so the comparison
    is paired.}
    \label{fig:unc-vs-mods}
\end{figure}

% --- Figure: qualitative --------------------------------------------------
% From eval_output/<run>/combo_grids/. Rows = configurations, columns = seeds.
\begin{figure}[t]
    \centering
    \includegraphics[width=\textwidth]{combo_grid.pdf}
    \caption{Samples for one case under decreasing modality availability.
    Rows are configurations ordered from most to least informative; columns
    are independent sampling seeds, followed by the ensemble and the
    uncertainty map.}
    \label{fig:combo-grid}
\end{figure}

% ---------------------------------------------------------------------------
% WRITING ORDER FOR THIS SUBSECTION once the numbers are in:
%
%   (a) Autoencoder ceiling -- one or two sentences on Table 1. State the
%       annotation reconstruction Dice explicitly; every later number must be
%       read against it.
%
%   (b) Accuracy across configurations -- Table 2, Dice and IoU columns.
%       Report the drop from four sequences to one. Note which single
%       sequence performs best; this is an empirical finding, not something
%       Eq. (dpi) predicts, since those configurations are not nested.
%
%   (c) The uncertainty result -- Table 2 U column and Fig. 5. State whether
%       U increases monotonically in the number of removed sequences, and
%       give the correlation. This is the claim the paper rests on.
%
%   (d) Failure cases -- configurations where accuracy is poor but
%       uncertainty stays low. These are the miscalibrated cases and reporting
%       them is more credible than omitting them.
%
%   (e) Qualitative -- Fig. 6. Point at the contrast between a thin boundary
%       band with four sequences and diffuse uncertainty with one.
% ---------------------------------------------------------------------------